% Source-based reconstruction: HoTT book official public-release, 19 June 2013.
% Source-derived material is attributed under the book's CC BY-SA 3.0 license.
\chapter{Higher Inductive Types and Cubical Type Theory}
\label{sec:p07-higher-inductive-types}

Higher inductive types add generators in identity types as well as generators
of points. Their induction principles belong to the type being generated;
its identity types retain ordinary identity elimination, not a new rule saying
that every path is a word in the displayed generators. This presentation
follows the official June 2013 public-release of the HoTT book
\cite[Chapter~6]{HoTT2013PublicReleaseFoundations}. It does not identify that source
with a later corrected build or with a computational cubical theory.

\section{Point and path constructors}

\begin{hbconvention}{def:p07-hit-standing}{Ambient theory and computation standing}
Work in intensional type theory with dependent products, dependent sums,
identity elimination, universes, the empty and unit types, and ordinary
inductive booleans, with judgmental function \(\eta\)-conversion as used
in the source's point computations. Paths compose from left to right: \(p\cdot q\)
traverses \(p\) and then \(q\). The notation \(=\) denotes an identity
type; \(\equiv_{\mathrm j}\) denotes judgmental equality. Function
extensionality is explicitly assumed for results asserting equalities of
functions and equivalences of mapping types, as in
\cref{def:p07-hott-funext,thm:p07-hott-function-paths}.
The interval argument below does not use that assumption, but uses the
ordinary judgmental function \(\eta\)-conversion.

The higher-inductive declarations below are added formation, introduction,
elimination and computation principles. They are not proofs of existence in
ordinary MLTT. In the 2013 presentation, computation at point constructors
is judgmental; computation at path and higher-path constructors is a specified
propositional equality. Large elimination is allowed where indicated; the
univalence example requires a larger universe containing the universe of its
boolean fiber. No normalization or canonicity theorem is asserted.
\end{hbconvention}

\begin{hbdefinition}{def:p07-hit-dependent-paths}{Paths over paths and paths over 2-paths}
For \(P:A\to\mathcal U\), \(p:x=_A y\), \(u:P(x)\), and
\(v:P(y)\), write
\[
 \operatorname{DPath}_P(p;u,v)
       :=\bigl(\operatorname{tr}_P(p,u)=_{P(y)}v\bigr).
\]
Here transport and dependent action have the types of
\cref{def:p07-hott-transport,thm:p07-hott-transport-laws}:
\(\operatorname{apd}_f(p):\operatorname{DPath}_P(p;f(x),f(y))\).
For \(\alpha:p=q\), put
\[
 \operatorname{tr}^2_P(\alpha,u)
 =\operatorname{ap}_{r\mapsto\operatorname{tr}_P(r,u)}(\alpha),\qquad
 \operatorname{DPath}^2_P(\alpha;h,k)
 :=\bigl(h=\operatorname{tr}^2_P(\alpha,u)\cdot k\bigr),
\]
where \(h:\operatorname{DPath}_P(p;u,v)\) and
\(k:\operatorname{DPath}_P(q;u,v)\). Both sides of the latter equality
are paths from \(\operatorname{tr}_P(p,u)\) to \(v\) in \(P(y)\).
This is the transport-based dependent 2-path convention of \S6.4.
\end{hbdefinition}

\begin{hblemma}{thm:p07-hit-path-calculus}{proved=annex:p07-hit-path-calculus}{The transport comparisons needed by eliminators}
In a constant family \(x\mapsto B\), there is a path
\(c(p,u):\operatorname{tr}_{x\mapsto B}(p,u)=u\), and
\[
 \operatorname{apd}_f(p)=c(p,f(x))\cdot\operatorname{ap}_f(p).
\]
For \(f,g:A\to B\), transport in their equality family satisfies
\[
 \operatorname{tr}_{x\mapsto f(x)=g(x)}(p,r)
       =\operatorname{ap}_f(p)^{-1}\cdot r\cdot\operatorname{ap}_g(p).
\]
For fixed \(b:B\), transport in \(x\mapsto x=b\) sends
\(r:x=b\) to \(p^{-1}\cdot r\). A dependent function has an action
\(\operatorname{apd}^2_f(\alpha)\) with the dependent 2-path type above.
These are propositional comparisons, not additional judgmental conversions.
Dependent paths compose by
\[
 h\star k:=\gamma_P(p,q,u)\cdot
       \operatorname{ap}_{\operatorname{tr}_P(q,-)}(h)\cdot k,
\]
where \(h:\operatorname{DPath}_P(p;u,v)\),
\(k:\operatorname{DPath}_P(q;v,w)\), and
\(\gamma_P(p,q,u):\operatorname{tr}_P(p\cdot q,u)=
\operatorname{tr}_P(q,\operatorname{tr}_P(p,u))\).
There is a comparison
\(\mu_f(p,q):\operatorname{apd}_f(p\cdot q)=
\operatorname{apd}_f(p)\star\operatorname{apd}_f(q)\).
\end{hblemma}

\begin{hbdefinition}{def:p07-hit-circle}{Circle constructors and the dependent eliminator}
The circle has a point \(\mathrm{base}:S^1\) and a path
\(\mathrm{loop}:\mathrm{base}=_{S^1}\mathrm{base}\).
For \(P:S^1\to\mathcal U\), data for induction are
\[
 b:P(\mathrm{base}),\qquad
 \ell:\operatorname{DPath}_P(\mathrm{loop};b,b).
\]
They give \(f:\prod_{x:S^1}P(x)\), with
\[
 f(\mathrm{base})\equiv_{\mathrm j}b,\qquad
 \beta_{\mathrm{loop}}:
             \operatorname{apd}_f(\mathrm{loop})=\ell.
\]
The second argument is not merely \(b=b\) in the base fiber: it includes
transport along the generating loop. The second computation clause is a path
in the dependent-path type. It does not make the generating loop judgmentally
reflexivity.
\index{circle!higher inductive type}
\end{hbdefinition}

\begin{hbtheorem}{thm:p07-hit-circle-recursion}{proved=annex:p07-hit-circle-recursion}{Circle recursion with correctly converted path computation}
For \(B:\mathcal U\), \(b:B\), and \(\ell:b=b\), circle induction
gives \(r:S^1\to B\) with
\(r(\mathrm{base})\equiv_{\mathrm j}b\) and a specified path
\(\operatorname{ap}_r(\mathrm{loop})=\ell\).
\end{hbtheorem}

\begin{hbtheorem}{thm:p07-hit-circle-universal}{proved=annex:p07-hit-circle-universal}{Uniqueness and the mapping property of the circle}
For \(f,g:S^1\to B\), a path \(p:f(\mathrm{base})=g(\mathrm{base})\)
and a coherence
\[
 \operatorname{ap}_f(\mathrm{loop})\cdot p
            =p\cdot\operatorname{ap}_g(\mathrm{loop})
\]
give a homotopy \(f\sim g\), hence \(f=g\) under function extensionality.
The evaluation map
\[
 (S^1\to B)\longrightarrow\sum_{b:B}(b=b),\qquad
 f\longmapsto(f(\mathrm{base}),\operatorname{ap}_f(\mathrm{loop}))
\]
has an explicit two-sided homotopy inverse given by recursion and is an
equivalence. Arbitrary point agreement without the displayed loop coherence
does not suffice.
\end{hbtheorem}

\begin{hbdefinition}{def:p07-hit-interval}{The higher-inductive interval}
The type \(I\) has points \(0_I,1_I:I\) and
\(\mathrm{seg}:0_I=1_I\). Induction into \(P:I\to\mathcal U\)
requires \(u:P(0_I)\), \(v:P(1_I)\), and
\(s:\operatorname{DPath}_P(\mathrm{seg};u,v)\); it gives a section
with judgmental endpoint computation and propositional \(\operatorname{apd}\)
computation at \(\mathrm{seg}\). In a constant family this gives recursion
from two endpoints and a path between their values, with \(\operatorname{ap}\)
computation. This \(I\) is a type with an identity constructor, not a
cubical dimension object with face formulas or a composition operation.
\end{hbdefinition}

\begin{hbtheorem}{thm:p07-hit-interval-contractible}{proved=annex:p07-hit-interval-contractible}{The interval is contractible}
There is a term \(\prod_{x:I}(x=1_I)\), with values \(\mathrm{seg}\)
at \(0_I\) and \(\mathrm{refl}_{1_I}\) at \(1_I\). Thus \(I\)
is contractible, although its endpoint declarations are not judgmentally the
same point.
\end{hbtheorem}

\begin{hbtheorem}{thm:p07-hit-interval-extensionality}{proved=annex:p07-hit-interval-extensionality}{An interval supplies a function-equality witness}
Without assuming function extensionality, if \(f,g:A\to B\) and
\(h:\prod_{a:A}(f(a)=g(a))\), interval recursion and judgmental function
\(\eta\) give a path \(f=g\). This is the displayed existence result of
\S6.3, not a claim here about judgmental equality or the full computation
laws of an inverse to \(\operatorname{happly}\).
\end{hbtheorem}

\begin{hbtheorem}{thm:p07-hit-circle-nontrivial}{proved=annex:p07-hit-circle-nontrivial}{The loop is nontrivial under univalence}
Assume a univalent universe containing booleans and circle elimination into
that universe. Then \(\mathrm{loop}\ne\mathrm{refl}_{\mathrm{base}}\).
The proof uses the boolean-swap equivalence as monodromy. No identification
\(\Omega(S^1)\simeq\mathbb Z\) is inferred from this single nontriviality
result.
\end{hbtheorem}

\section{Suspensions and homotopy pushouts}

\begin{hbdefinition}{def:p07-hit-suspension}{Suspension constructors and eliminators}
For \(A:\mathcal U\), its suspension \(\Sigma A\) has points
\(\mathrm{north},\mathrm{south}\) and
\(\mathrm{merid}:\prod_{a:A}(\mathrm{north}=\mathrm{south})\).
Here \(\Sigma A\) is the source's suspension notation, distinct from a
dependent-sum expression \(\sum_{x:A}P(x)\).
For \(P:\Sigma A\to\mathcal U\), induction data are
\[
 n:P(\mathrm{north}),\quad s:P(\mathrm{south}),\quad
 m:\prod_{a:A}\operatorname{DPath}_P(\mathrm{merid}(a);n,s).
\]
The section computes judgmentally at the two poles and propositionally by
\(\operatorname{apd}\) on each meridian. Recursion into \(B\) takes
\(n,s:B\) and \(m:A\to(n=s)\), with the corresponding \(\operatorname{ap}\)
computation. The dependence on \(a\) also transmits paths and higher paths
of \(A\); it is not a collection indexed only by the set of components.
\end{hbdefinition}

\begin{hbtheorem}{thm:p07-hit-suspension-circle}{proved=annex:p07-hit-suspension-circle}{Suspending booleans gives the circle}
There is an equivalence \(\Sigma\mathbf2\simeq S^1\), with forward
map sending both poles to \(\mathrm{base}\), the false meridian to
\(\mathrm{loop}\), and the true meridian to reflexivity.
Its inverse sends \(\mathrm{base}\) to \(\mathrm{north}\) and
\(\mathrm{loop}\) to
\(\mathrm{merid}(\mathrm{false})\cdot\mathrm{merid}(\mathrm{true})^{-1}\).
The equivalence requires the two meridian coherences, not only the pole values.
\end{hbtheorem}

\begin{hbdefinition}{def:p07-hit-spheres}{Spheres and the typed 2-dimensional clause}
The suspension presentation is \(S^{-1}:=\mathbf0\),
\(S^0:=\mathbf2\), and \(S^{n+1}:=\Sigma S^n\) for \(n\ge0\).
The direct 2-sphere presentation has \(\mathrm{base}:S^2\) and
\(\mathrm{surf}:\mathrm{refl}_{\mathrm{base}}=
\mathrm{refl}_{\mathrm{base}}\) in its loop type.
For \(P:S^2\to\mathcal U\), induction takes \(b:P(\mathrm{base})\)
and \(s:\operatorname{DPath}^2_P(\mathrm{surf};\mathrm{refl}_b,
\mathrm{refl}_b)\), computing judgmentally at the point and
propositionally by \(\operatorname{apd}^2\) at \(\mathrm{surf}\).
The direct and iterated-suspension presentations are not identified without
the required comparison proof. Constructor dimension alone does not bound
the dimensions of the identity types.
\end{hbdefinition}

\begin{hbdefinition}{def:p07-hit-pushout}{The dependent pushout principle}
For a span \(A\xleftarrow{u}C\xrightarrow{v}B\), the pushout \(W\)
has \(\mathrm{inl}:A\to W\), \(\mathrm{inr}:B\to W\), and
\(\mathrm{glue}(c):\mathrm{inl}(u(c))=\mathrm{inr}(v(c))\).
For \(P:W\to\mathcal U\), induction requires
\[
 d_L:\prod_{a:A}P(\mathrm{inl}(a)),\quad
 d_R:\prod_{b:B}P(\mathrm{inr}(b)),\quad
 e:\prod_{c:C}\operatorname{DPath}_P
      (\mathrm{glue}(c);d_L(u(c)),d_R(v(c))).
\]
The resulting section computes judgmentally on \(\mathrm{inl}\) and
\(\mathrm{inr}\), and propositionally by \(\operatorname{apd}\) on
\(\mathrm{glue}\). Recursion is the constant-family case.
\end{hbdefinition}

\begin{hbtheorem}{thm:p07-hit-pushout-universal}{proved=annex:p07-hit-pushout-universal}{Maps out of a pushout are coherent cocones}
Under function extensionality, evaluation on the constructors is an equivalence
\[
 (W\to E)\simeq
 \sum_{i:A\to E}\sum_{j:B\to E}
                \prod_{c:C}(i(u(c))=j(v(c))).
\]
The homotopy pushout of \(\mathbf1\leftarrow A\to\mathbf1\) is a
suspension presentation. In particular, a homotopy pushout of sets need not
be a set: with middle type \(\mathbf2\), it is the preceding circle
presentation, which is not a set under the stated univalence assumption.
\end{hbtheorem}
\section{Truncations and set quotients}

\begin{hbdefinition}{def:p07-hit-truncations}{The recursive propositional and set-truncation declarations}
The propositional truncation \(\|A\|\) has point constructor
\(|-|:A\to\|A\|\) and path constructor
\(\mathrm{squash}:\prod_{x,y:\|A\|}(x=y)\). For a family
\(P:\|A\|\to\mathcal U\), its raw dependent elimination data are
\(g:\prod_{a:A}P(|a|)\) and, for every \(x,y,u:P(x),v:P(y)\),
a lift \(\operatorname{DPath}_P(\mathrm{squash}(x,y);u,v)\).
The section satisfies \(f(|a|)\equiv_{\mathrm j}g(a)\) and the
specified propositional action on \(\mathrm{squash}\).
In particular a proposition-valued family supplies all those lifts: both
transported and target elements lie in the same proposition.

The set truncation \(\|A\|_0\) has \(|-|_0:A\to\|A\|_0\) and
\[
 \mathrm{trunc}:\prod_{x,y:\|A\|_0}\prod_{p,q:x=y}(p=q).
\]
For \(P:\|A\|_0\to\mathcal U\), raw elimination takes
\(g:\prod_{a:A}P(|a|_0)\) and, for each such \(x,y,p,q\), each
\(u:P(x),v:P(y)\), and lifts \(h:\operatorname{DPath}_P(p;u,v)\),
\(k:\operatorname{DPath}_P(q;u,v)\), a term
\(\operatorname{DPath}^2_P(\mathrm{trunc}(x,y,p,q);h,k)\).
It has judgmental point computation and propositional computation on the
2-path constructor. The 2013 text notes that identity-type inputs of this
direct declaration fall outside its simpler strict-positive-input syntax;
the hub-and-spoke construction below is its separate systematic alternative.
\index{truncation!propositional}\index{truncation!set}
\end{hbdefinition}

\begin{hbtheorem}{thm:p07-hit-truncation-elimination}{proved=annex:p07-hit-truncation-elimination}{Restricted dependent elimination and universal properties}
Propositional truncation eliminates into proposition-valued families from
the point data alone. Set truncation eliminates into set-valued families
from the point data alone, because their dependent 2-path coherence exists.
Under function extensionality, precomposition with the point constructors
gives equivalences
\[
 (\|A\|\to B)\simeq(A\to B)\quad\text{if }B\text{ is a proposition},
 \qquad
 (\|A\|_0\to B)\simeq(A\to B)\quad\text{if }B\text{ is a set}.
\]
The restrictions on \(B\) are part of the conclusions, not dispensable
annotations.
\end{hbtheorem}

\begin{hbtheorem}{thm:p07-hit-truncation-boundaries}{proved=annex:p07-hit-truncation-boundaries}{Empty cases and failure of unrestricted witness extraction}
\(\|\mathbf0\|\simeq\mathbf0\), while \(\|\mathbf1\|\) and
\(\|\mathbf2\|\) are contractible. There is no map
\(e:\|\mathbf2\|\to\mathbf2\) with \(e(|b|)=b\) for both boolean
points. Thus the proposition merely asserting an inhabitant is not a
choice of an arbitrary witness in a nonpropositional target.
\end{hbtheorem}

\begin{hbdefinition}{def:p07-hit-set-quotient}{Set-quotient constructors and dependent elimination}
Let \(A\) be a set and \(R:A\to A\to\mathrm{Prop}\) a mere relation,
not necessarily an equivalence relation. The set-quotient \(Q=A/R\) has
\[
 q:A\to Q,\qquad
 \mathrm{eq}:\prod_{a,b:A}(R(a,b)\to(q(a)=q(b))),\qquad
 \mathrm{trunc}_Q:\prod_{x,y:Q}\prod_{p,s:x=y}(p=s).
\]
It is a set by its final constructor. Induction into a set-valued family
\(P:Q\to\mathcal U\) takes \(d:\prod_aP(q(a))\) and
\[
 e:\prod_{a,b:A}\prod_{r:R(a,b)}
             \operatorname{DPath}_P(\mathrm{eq}(r);d(a),d(b)).
\]
The 2-path lifts are supplied as for set truncation. Point computation is
judgmental and \(\mathrm{eq}\)-computation is propositional.
In a constant set \(B\), these data are a map \(g:A\to B\) and
\(\rho:\prod_{a,b}R(a,b)\to(g(a)=g(b))\).
\end{hbdefinition}

\begin{hbtheorem}{thm:p07-hit-quotient-universal}{proved=annex:p07-hit-quotient-universal}{Mere surjectivity and the set-quotient mapping property}
The projection is merely surjective:
\[
 \prod_{x:Q}\left\|\sum_{a:A}(q(a)=x)\right\|.
\]
For any set \(B\), evaluation is an equivalence
\[
 (Q\to B)\simeq
 \sum_{g:A\to B}\prod_{a,b:A}(R(a,b)\to(g(a)=g(b))).
\]
This gives no global representative selector. For arbitrary \(R\) it also
does not give \(R(a,b)\simeq(q(a)=q(b))\): the identity type includes
reflexivity, symmetry and composition, whereas \(R\) need not.
\end{hbtheorem}

\begin{hbtheorem}{thm:p07-hit-quotient-example}{proved=annex:p07-hit-quotient-example}{A directed relation can generate a one-class quotient}
Let \(A=\mathbf2\) and let \(R(\mathrm{false},\mathrm{true})=\mathbf1\)
with all three other boolean pairs assigned \(\mathbf0\).
Then \(A/R\) is contractible, although \(R\) is neither reflexive nor
symmetric. In particular, \(q(\mathrm{false})=q(\mathrm{false})\)
is inhabited while \(R(\mathrm{false},\mathrm{false})\) is empty.
\end{hbtheorem}

\begin{hbdefinition}{def:p07-hit-hub-truncation}{The original general hub-and-spoke declaration}
For \(n\ge-1\), the general truncation construction of \S7.3 has
\(\eta_n:A\to\|A\|_n\), and for every
\(r:S^{n+1}\to\|A\|_n\) a point \(\mathrm{hub}(r)\) and paths
\(\mathrm{spoke}_r(x):r(x)=\mathrm{hub}(r)\).
For \(P:\|A\|_n\to\mathcal U\), its full raw dependent data are
\[
 g:\prod_aP(\eta_n(a)),\quad
 H:\prod_r\prod_{r':\prod_xP(r(x))}P(\mathrm{hub}(r)),
 \quad
 S:\prod_r\prod_{r'}\prod_x
       \operatorname{DPath}_P(\mathrm{spoke}_r(x);r'(x),H(r,r')).
\]
They supply a section \(f\), with judgmental point computation
\(f(\eta_n(a))\equiv_{\mathrm j}g(a)\) and
\(f(\mathrm{hub}(r))\equiv_{\mathrm j}H(r,f\circ r)\), and
propositional computation
\(\operatorname{apd}_f(\mathrm{spoke}_r(x))=S(r,f\circ r,x)\).
The type \(\|A\|_{-2}\) is instead defined as \(\mathbf1\), not by
this scheme. Removing the hub and connecting each \(r(x)\) directly to
\(r(\mathrm{base})\) leaves an uncontrolled basepoint loop and is not
silently substituted for the source construction.
\end{hbdefinition}

\begin{hbtheorem}{thm:p07-hit-hub-contraction}{proved=annex:p07-hit-hub-contraction}{The exact based-mapping contraction supplied by spokes}
For \(X=\|A\|_n\), any pointed \(S^{n+1}\) with basepoint \(s_0\),
and \(b:X\), the type
\(\sum_{r:S^{n+1}\to X}(r(s_0)=b)\) is contractible.
For \((r,p)\), the pointwise homotopy to the constant map at \(b\) is
\[
 x\longmapsto\mathrm{spoke}_r(x)\cdot
                 \mathrm{spoke}_r(s_0)^{-1}\cdot p.
\]
The proof includes transport of the basepoint path, not only equality of the
underlying maps.
\end{hbtheorem}

\begin{hbtheorem}{thm:p07-hit-hub-dependent}{proved=annex:p07-hit-hub-dependent}{Dependent induction from explicitly given fiber contractions}
Suppose, for each fiber \(P(z)\), a function is given taking every
\(t:S^{n+1}\to P(z)\) to a point \(u:P(z)\) and a homotopy
\(\prod_x(t(x)=u)\). Then point data
\(g:\prod_aP(\eta_n(a))\) extend to a section of \(P\) with its
judgmental point computation. The supplied fiber-contraction function is a
specific additional hypothesis; it is not an assertion that such a function
has already been constructed for all \(n\)-types.
\end{hbtheorem}

\begin{proofboundary}{def:p07-hit-general-boundaries}{General truncation, homotopy groups, and source variants}
The full conversion of the preceding based-sphere mapping contraction into
\(n\)-truncatedness uses the \(\Omega^{n+1}\)-characterization of
\(n\)-types and the sphere--loop adjunction, including its based-path
coherences. Likewise deriving the restricted general \(n\)-type eliminator
requires the corresponding chosen contractions of sphere maps into each
fiber. Those dependencies must be supplied before calling this a completed
general-truncation proof; the explicit low-dimensional truncation proofs above
do not replace them.

The calculation \(\Omega(S^1)\simeq\mathbb Z\) needs the actual code
family, winding/transport calculation, both encode--decode composites and their
coherence. The circle constructors or nontrivial-loop argument do not give
that calculation. Higher sphere homotopy groups are not asserted here.

The original source also has direct cell-complex and hub presentations,
impredicative quotients, free/presented algebras, and the flattening lemma
\cite[\S\S6.6--6.7, 6.10--6.13]{HoTT2013PublicReleaseFoundations}. They retain their
own constructions and hypotheses. The exact coequalizer flattening theorem,
its total-space constructors, and both inverse homotopies are supplied in
Section~\ref{sec:p07-hott-identity-systems-and-flattening}.
In particular its impredicative
quotient of equivalence classes uses propositional univalence and raises a
universe unless resizing is supplied. It is not identified with the HIT
quotient by a title match. The 2013 text gives no complete general syntactic
criterion or MLTT-only existence theorem for every proposed HIT: ordered
constructors, positivity and naturality are conditions to be checked, not a
license to add arbitrary polymorphic endpoints. Computational cubical
composition and judgmental path computation belong to their separate source
presentation.
\end{proofboundary}
